\chapter{The theoretical part}
\label{ch:method}

This chapter describes the design and methodology that produces every number in the diploma project.  All code lives in \file{src/} and \file{experiments/}; YAML configs in \file{experiments/configs/} are the canonical record of every run.

\section{Unified detector interface}

\file{src/models/base.py} defines an abstract \code{AnomalyDetector} with three methods:
\begin{itemize}
  \item \code{fit(X\_train, X\_val)} --- trained on normal data only.
  \item \code{score(X)}$\rightarrow$\code{(N,)} --- per-sample anomaly score, higher $=$ more anomalous.
  \item \code{attribute(X)}$\rightarrow$\code{(N, F)} --- per-feature contribution; default raises, overridden by AE family and TranAD.
\end{itemize}

All six detectors implement this interface.  The runner (\file{experiments/run.py}) treats them uniformly, so cross-dataset study, LOPO, and attribution use the same code path as same-dataset detection.

\section{Data loaders (\file{src/data\_loader.py})}

Both loaders return a \code{DatasetBundle} with \code{features}, \code{labels}, \code{attack\_ids}, \code{split}, \code{name}, \code{timestamps}, and \code{metadata}.

\paragraph{HAI 21.03.}  Consolidates per-process \code{attack\_P\{n\}} flags into a single \code{label} and keeps the first-fired flag as the per-row \code{attack\_id}.  Train/val split is random over the normal \code{train*.csv.gz} files; test is the full attack-mixed set.

\paragraph{Morris gas-pipeline.}  ARFF parser with a \code{pressure measurement} sentinel fix --- values near \code{float32}-max ($\sim 3.4\!\times\!10^{38}$) are replaced with 0 because MinMax scaling and the schema-align mean were otherwise dominated by them.

\paragraph{Leak assertion.}  \code{DatasetBundle.assert\_no\_attack\_in\_train\_val} runs before any downstream step.  Attack rows in train/val are the most common source of inflated published numbers; we fail loudly.

\section{Preprocessing (\file{src/preprocessing.py})}

\begin{itemize}
  \item \code{scale\_bundle(bundle)}$\rightarrow$\code{ScaledArrays} fits \code{MinMaxScaler} on \textbf{train only}, transforms all splits.  Assertion triggers before fit.
  \item \code{make\_windows(X, window, stride)} produces $(N, W, F)$ sliding windows for windowed models.
  \item \code{window\_labels(y, window, stride)} tags a window as 1 iff any frame inside is anomalous.
  \item \code{percentile\_threshold(val\_scores, 99.0)} picks the operating point from validation scores (never from test).
\end{itemize}

Each step has an exact definition that the rest of the apparatus depends on.  \emph{Scaling} is per-feature min--max with the statistics estimated on the training-normal split only,
\begin{equation}
  \tilde{x}_j = \frac{x_j - \min_{t \in \mathrm{train}} x_{t,j}}{\max_{t \in \mathrm{train}} x_{t,j} - \min_{t \in \mathrm{train}} x_{t,j}},
  \label{eq:minmax}
\end{equation}
with clipping to $[0,1]$ at transform time so that test values outside the training range do not produce out-of-bound inputs to the sigmoid-output autoencoders.  Fitting the scaler on normal data only --- rather than on the full series --- is the single most common source of inflated published numbers in ICS AD, because attack-range statistics leak the anomaly into the very normalisation the model trains against; the loader asserts the absence of attack rows in train and validation \emph{before}~\eqref{eq:minmax} is fit.  \emph{Windowing} maps a scaled matrix $X \in \mathbb{R}^{T \times F}$ to a tensor of $n = \lfloor (T - w)/s \rfloor + 1$ overlapping windows of length $w$ at stride $s$, and a window inherits the positive label iff any frame inside it is anomalous, $y^{(w)}_i = \max_{\,0 \leq \delta < w} y_{\,i s + \delta}$.  \emph{Thresholding} selects the operating point as the $p$-th percentile of the anomaly scores on the clean validation fold,
\begin{equation}
  \tau = \mathrm{Percentile}_{p}\bigl(\{\, s(x) : x \in \mathrm{val} \,\}\bigr), \qquad p = 99,
  \label{eq:threshold}
\end{equation}
so that, by construction, roughly $1\%$ of clean data is flagged and no test label ever participates in threshold selection.  Equation~\eqref{eq:threshold} is the source-calibration rule; the target-calibration variant of Section~\ref{ch:results-transfer} substitutes the target domain's clean validation scores for the source domain's.

\section{Evaluation (\file{src/evaluation/})}

\begin{itemize}
  \item \file{pointwise.py} --- precision, recall, F1, ROC-AUC, PR-AUC.
  \item \file{point\_adjust.py} --- Xu-2018 PA-F1.
  \item \file{etapr.py} --- Hwang-2022 event-weighted precision/recall, ported from the reference implementation and verified against their published numbers on a known model.
\end{itemize}

All three return a tiny metric class with an \code{as\_dict()} serializer; the runner concatenates one row per metric into \file{results/metrics/summary.parquet}.

\section{Cross-testbed transfer (\file{src/transfer/schema\_align.py})}

The method is:
\begin{enumerate}
  \item Hand-tag every HAI and Morris feature with a canonical type (\code{pressure}, \code{pump\_state}, \code{setpoint}, \code{system\_state}, \code{valve\_position}, \code{control\_signal}; plus \code{unknown}/\code{comm\_metadata} which are excluded).  Tags live in \file{data/feature\_types.yaml}.
  \item For each dataset, aggregate same-typed features into a single type column (default: mean; \code{max} for binary-state types via the \code{aggregations} block).  This produces a 6-column type-vector.
  \item Project both datasets to the intersection of types.  Train on source type-vector; score on target type-vector.
\end{enumerate}

Formally, the type tags induce a partition of each dataset's raw columns into type groups.  For dataset $d$ with feature set $\mathcal{F}_d$ and a tagging map $\mathrm{type}_d : \mathcal{F}_d \to \mathcal{T}$, the projection of a raw row $x$ onto the shared type space is the aggregation
\begin{equation}
  \Pi_d(x)_t = \operatorname{agg}_t\bigl(\{\, x_f : f \in \mathcal{F}_d,\; \mathrm{type}_d(f) = t \,\}\bigr),
  \qquad t \in \mathcal{T}_\cap,
  \label{eq:projection}
\end{equation}
where $\mathcal{T}_\cap = \mathrm{types}(\mathcal{F}_{\mathrm{src}}) \cap \mathrm{types}(\mathcal{F}_{\mathrm{tgt}}) \setminus \{\texttt{unknown}, \texttt{comm\_metadata}\}$ is the set of canonical types present in both datasets after exclusions, and $\operatorname{agg}_t$ is the per-type aggregator (mean by default; \code{max} for binary-state types such as \texttt{pump\_state} and \texttt{valve\_position}).  Because $\Pi_{\mathrm{src}}$ and $\Pi_{\mathrm{tgt}}$ both map into $\mathbb{R}^{|\mathcal{T}_\cap|}$, a detector trained on $\Pi_{\mathrm{src}}(X_{\mathrm{src}})$ can be evaluated directly on $\Pi_{\mathrm{tgt}}(X_{\mathrm{tgt}})$ with no shape mismatch.  The projection in~\eqref{eq:projection} is deliberately lossy: it discards per-variable identity in exchange for a cross-testbed input shape every detector can consume, and Section~\ref{ch:results-transfer} measures exactly how much detection capability that trade costs.

Two calibration regimes are implemented in \file{experiments/run\_transfer.py}:
\begin{itemize}
  \item \code{val\_percentile} --- threshold from 99th percentile of \emph{source} validation scores.
  \item \code{target\_val\_percentile} --- threshold from 99th percentile of \emph{target}-normal validation scores.  Decouples representation transfer from operating-point transfer.
\end{itemize}
The distinction matters because~\eqref{eq:projection} transfers the learned \emph{representation} (the model weights and the scoring rule) but says nothing about the \emph{operating point}: a threshold calibrated on the source-score distribution need not sit anywhere sensible on the target-score distribution.  Reporting both regimes is what lets Section~\ref{ch:results-transfer} separate a representation that transfers from a threshold that does not.

\section{Within-HAI LOPO (\file{src/transfer/lopo.py})}

\code{drop\_process\_features(bundle, "P2")} returns a new bundle with all \code{P2\_*} columns removed.  \code{per\_attack\_process\_f1(y\_true, y\_pred, attack\_ids)} computes F1 restricted to each process's attack subset (union with normal rows).  \file{experiments/run\_lopo.py} runs the full pipeline with a chosen process held out and emits \code{f1\_attack\_P\{1..4\}} columns in \file{summary.parquet}.

\section{Attribution evaluation (\file{src/attribution/evaluation.py})}

\begin{itemize}
  \item \code{feature\_to\_process(feature\_names)} maps HAI feature names to their \code{P\{n\}} prefix.
  \item \code{process\_precision\_at\_k(scores\_F, feature\_processes, attacked\_process, k)} --- fraction of top-k features that belong to the attacked process.
  \item \code{precision\_at\_k\_by\_attack} aggregates across a batch of attack windows, keyed by attacked process.
\end{itemize}

Formally, for an attack sample with per-feature attribution vector $a \in \mathbb{R}^{F}$, let $\mathrm{Top}_k(a)$ be the index set of the $k$ largest entries and let $\mathcal{P}^\star$ be the process under attack.  Precision@$k$ is the fraction of the top-$k$ attributed features that belong to that process,
\begin{equation}
  \mathrm{p}@k = \frac{\bigl|\{\, j \in \mathrm{Top}_k(a) : \mathrm{proc}(j) = \mathcal{P}^\star \,\}\bigr|}{k},
  \label{eq:patk}
\end{equation}
where $\mathrm{proc}(j)$ is the process prefix of feature $j$ (the \code{P\{n\}} tag).  The random baseline for~\eqref{eq:patk} --- the expected precision of a uniform-random attribution --- equals the share of features belonging to the attacked process,
\begin{equation}
  \mathrm{p}@k_{\mathrm{rand}} = \frac{|\{\, j : \mathrm{proc}(j) = \mathcal{P}^\star \,\}|}{F},
\end{equation}
independent of $k$.  We report it alongside the model numbers so that \emph{lift} --- the ratio $\mathrm{p}@k / \mathrm{p}@k_{\mathrm{rand}}$ --- is visible at a glance; a model that localizes no better than chance sits at lift~$1$.

\subsection{TranAD attention rollout}
\label{sec:tranad-attn}

\code{TranADModel.attribute\_attention(X)} hooks the encoder's self-attention, averages over heads and over output positions to get a per-input-timestep weight, and weights the per-$(t,f)$ squared reconstruction error by that weight.  Single-layer rollout degenerates to the attention matrix.  This is a literal reading of ``attention rollout as attribution'' for the TranAD architecture.

\section{Experiment infrastructure}

Four drivers, all config-YAML-driven:
\begin{itemize}
  \item \file{experiments/run.py} --- same-dataset detection.
  \item \file{experiments/run\_transfer.py} --- cross-testbed transfer under both calibration regimes.
  \item \file{experiments/run\_lopo.py} --- within-HAI leave-one-process-out.
  \item \file{experiments/run\_attribution.py} --- per-feature attribution with \code{--method \{reconstruction,attention\}} flag.
\end{itemize}

Results land in \file{results/metrics/summary.parquet} (detection) and \file{attribution.parquet} (attribution).  Seeds are set via \code{src.utils.set\_seed} at the top of every runner; config hashes are written into every result row so a saved number can be traced back to its exact config.

\section{Tests}

Mandatory for \file{evaluation/} and \file{data\_loader.py}; smoke coverage for everything else (CLAUDE.md policy).  Current suites: \file{test\_data\_loader.py}, \file{test\_preprocessing.py}, \file{test\_evaluation.py}, \file{test\_transfer.py}, \file{test\_attribution.py}, \file{test\_models\_smoke.py}.

The evaluation tests are hand-computed: a small synthetic label/prediction pair whose point-wise, point-adjust, and eTaPR scores are worked out on paper and asserted exactly, so that any future refactor that silently changes a metric definition fails loudly.  This matters for the thesis claim --- ranking changes across metrics (Section~\ref{ch:results-detection}) are only credible if the metrics themselves are pinned by tests rather than reimplemented per notebook.

\section{Computational complexity}
\label{sec:complexity}

The detectors span three orders of magnitude in training cost, which is the practical reason the benchmark is feasible on a single 4\,GB GPU.  Table~\ref{tab:complexity} summarises the asymptotic training and scoring cost of each model, with $n$ the number of training samples, $F$ the feature count, $w$ the window length, $d$ a representative hidden width, and $E$ the epoch count.  The classical baselines are effectively free to score; the windowed deep models dominate the wall-clock budget, which is why the reproducibility checks (\S\ref{sec:tranad-repro}) run on a $30\mathrm{k}$-window subsample rather than the full $\sim 1.5\mathrm{M}$ windows.

\begin{table}[h]
\centering
\caption{Asymptotic training and per-sample scoring complexity of the six detectors.}
\label{tab:complexity}
\small
\begin{tabular}{lll}
\toprule
Detector & Training & Scoring (per sample) \\
\midrule
Isolation Forest & $O(t\, \psi \log \psi)$ ($t$ trees, $\psi$ subsample) & $O(t \log \psi)$ \\
One-Class SVM    & $O(n_{\mathrm{sv}}^2 F)$ to $O(n^3)$ worst case & $O(n_{\mathrm{sv}} F)$ \\
Dense AE         & $O(E\, n\, F d)$ & $O(F d)$ \\
LSTM-AE          & $O(E\, n\, w\, d^2)$ & $O(w\, d^2)$ \\
USAD             & $O(E\, n\, w F d)$ & $O(w F d)$ \\
TranAD           & $O(E\, n\, w^2 d)$ (attention) & $O(w^2 d)$ \\
\bottomrule
\end{tabular}
\end{table}

The $w^2$ term in TranAD is the self-attention cost over the $60$-step window; it is the reason TranAD is the most expensive model per epoch despite its single encoder layer, and it bounds how deep the encoder can grow on the available hardware.

\section{Experiment lifecycle}
\label{sec:lifecycle}

Every reported number flows through the same five-stage pipeline regardless of which of the four drivers invokes it: load, scale (with the leak guard), fit on normal data, threshold on clean validation, and score the labelled test split under all three metrics.  Listing~\ref{lst:runner} sketches the same-dataset driver; the transfer, LOPO, and attribution drivers differ only in which bundle transformation (\eqref{eq:projection} or process-dropping) is applied between loading and scaling, and in which threshold distribution feeds~\eqref{eq:threshold}.

\begin{lstlisting}[language=Python, caption={Same-dataset detection driver (abridged from \file{experiments/run.py}).}, label={lst:runner}]
def run(cfg):
    set_seed(cfg.seed)                      # reproducibility: one seed, recorded
    bundle = load_dataset(cfg.dataset)      # unified DatasetBundle
    scaled = scale_bundle(bundle)           # MinMax on TRAIN only; asserts no leak

    X_tr, X_val, X_te = scaled.X_train, scaled.X_val, scaled.X_test
    if cfg.windowed:                        # LSTM-AE / USAD / TranAD
        X_tr  = make_windows(X_tr,  cfg.window)
        X_val = make_windows(X_val, cfg.window)
        X_te  = make_windows(X_te,  cfg.window)
        y_te  = window_labels(scaled.y_test, cfg.window)
    else:
        y_te  = scaled.y_test

    model = build_model(cfg)                # dispatch on cfg.model name
    model.fit(X_tr, X_val)                  # normal-only fit

    val_scores = model.score(X_val)
    tau = percentile_threshold(val_scores, 99.0)   # never sees test labels
    test_scores = model.score(X_te)

    row = {"config_hash": cfg.hash, "seed": cfg.seed}
    for metric in (pointwise_metrics, point_adjust_f1, etapr_f1):
        row.update(metric(y_te, test_scores, tau).as_dict())
    append_parquet("results/metrics/summary.parquet", row)
\end{lstlisting}

The single \code{set\_seed} call, the train-only scaler fit, the validation-derived threshold, and the recorded \code{config\_hash} are the four invariants that make every cell of every results table in Chapter~\ref{ch:implementation} reproducible from a committed configuration.
